\documentclass[11pt,a4paper]{article}
\usepackage[margin=25mm]{geometry}
\usepackage{amsmath,amssymb,amsthm,hyperref}
\newtheorem{lemma}{Lemma}\newtheorem{proposition}[lemma]{Proposition}
\newcommand{\Sz}{\operatorname{Sz}}\newcommand{\gen}[1]{\langle#1\rangle}
\title{The remaining subfield obstruction for Suzuki rotation markings}
\author{Human Compression project: private F003 continuation}\date{}
\begin{document}\maketitle
This is a bounded continuation of the accepted proper-face argument, not
a proof of Spiga's unrestricted Theorem 1.1. Its conclusion identifies
the outstanding case exactly. The historical assertion and attribution
remain Spiga's; there is no new conjecture, original, formal or numerical
credit. The root reviewer independently checked the mixed-face argument
before this draft; the final manuscript has not yet had distinct review.

\paragraph{Paid and retained premises.}
We use the entire two-page \emph{Proper face groups in a Suzuki rotation
marking}, in \texttt{corpus106\_F003\_proper\_faces\_v1}, including its
root-authored matrix obstruction, odd-order endpoint lemma and retained
rotation-polytope interfaces. We additionally state the classical ovoid
interfaces needed here. For $S=\Sz(q)$ and a canonical subfield group
$H=\Sz(q_0)$, $q=q_0^m$ with $m$ odd, there is the canonical subovoid
$\Omega_H\subseteq\Omega$. Both actions are doubly transitive; their
point root groups $U_Q$ and $U_Q\cap H$ act regularly on
$\Omega\setminus\{Q\}$ and $\Omega_H\setminus\{Q\}$ respectively.
Every nonidentity even-order element has its unique fixed point; split
odd torus elements have exactly two fixed points and nonsplit odd torus
elements none. The odd tori have orders $q-1$ and
$q\pm\sqrt{2q}+1$. The root-group and subovoid inclusions agree in the
standard matrix coordinates. These are explicit classical Suzuki
structure premises, also used in the accepted Topic D and L015 papers;
their proofs are not replaced by this note.

\begin{lemma}[Fixed-point descent]\label{fixed}
If $1\ne h\in H$ fixes $Q\in\Omega$, then $Q\in\Omega_H$.
\end{lemma}
\begin{proof}
If $h$ has even order, its unique fixed point in the smaller ovoid is
also an ambient fixed point, and ambient uniqueness proves the claim.
If $h$ is in a split odd torus of $H$, its two subovoid fixed points
are its only ambient fixed points, by the retained Zassenhaus property.
For the remaining case, $|h|>1$ divides $q_0^2+1$. An odd-order element
fixing an ambient point has order dividing $q-1$: in a point stabilizer
the root group is a $2$-group and its odd complement has order $q-1$.
But
\[
 \gcd(q_0^2+1,q_0^m-1)=1 \qquad(m\ \hbox{odd}).
\]
Indeed a common divisor is odd, and modulo it one has $q_0^2=-1$ and
$q_0^m=1$, so $1=q_0^{2m}=(-1)^m=-1$. The divisor divides two and
is therefore one. Thus a nonsplit odd torus element cannot fix an
ambient point. This exhausts the element types.
\end{proof}

\begin{lemma}[One root-group move descends]\label{move}
Let $Q\in\Omega_H$ and $x,y\in\Omega_H\setminus\{Q\}$. If
$v\in U_Q$ sends $x$ to $y$, then $v\in H$.
\end{lemma}
\begin{proof}
Regularity of $U_Q\cap H$ supplies an element $v_0$ with the same
effect on $x$. Regularity of the whole $U_Q$ makes that element unique,
so $v=v_0$.
\end{proof}

\begin{proposition}[Mixed faces are impossible]\label{mixed}
Suppose $S=\gen{a,b,c}$, the products $ab,bc,abc$ are involutions,
$|a|,|c|>2$ and $b^2\ne1$. If $F=\gen{a,b}$ is contained in a
proper Suzuki subfield group and $E=\gen{b,c}$ fixes an ovoid point,
then a contradiction follows. The dual assertion holds as well.
\end{proposition}
\begin{proof}
Write $F\le H=\Sz(q_0)$ and let $E$ fix $Q$. Since $1\ne b\in H$
fixes $Q$, Lemma~\ref{fixed} gives $Q\in\Omega_H$. The involution
$v=bc$ belongs to the root group $U_Q$ of the point stabilizer. Also
$av=abc$ is an involution, whence $vav=a^{-1}$. Thus $v$, which fixes
$Q$, sends $aQ$ to $a^{-1}Q$. If $aQ=Q$, all three generators fix
$Q$, contradicting generation of $S$. Otherwise both points lie in
$\Omega_H\setminus\{Q\}$, and Lemma~\ref{move} implies $v\in H$.
Now $c=b^{-1}v\in H$, again contradicting generation. Duality
interchanges the two end-face groups and yields the other case.
\end{proof}

\begin{proposition}[Exact remaining case]\label{remaining}
Any rank-four chiral-polytope rotation marking of $\Sz(2^N)$ has
both proper face groups equal to conjugate subfield Suzuki groups
$\Sz(2^d)$ and $\Sz(2^e)$, where $d,e\mid N$ and
$1<d,e<N$. In particular $N$ must be composite.
\end{proposition}
\begin{proof}
The paid proper-face paper establishes $|a|,|c|>2$, $b^2\ne1$,
properness of $F,E$, and excludes torus-normalizer containment of either
face and simultaneous Borel containment. If one face lies in a Borel,
the maximal-subgroup classification places the other in either a Borel,
a torus normalizer or a proper subfield group. The first two alternatives
were excluded, and Proposition~\ref{mixed} excludes the third. Thus
neither face lies in an ambient Borel or torus normalizer.

Choose a smallest subfield Suzuki group containing $F$, available from
the retained maximal-subgroup list. If $F$ is proper in it, apply that
same list within the smaller Suzuki group. Its Borels lie in ambient
Borels. Its odd-torus normalizers lie in ambient odd-torus normalizers:
they normalize their nontrivial characteristic odd cyclic torus, hence
its unique ambient torus centralizer. A further subfield contradicts
minimality. The exceptional bottom group $\Sz(2)=C_5\rtimes C_4$
also normalizes its characteristic $C_5$ and therefore an ambient torus.
All proper alternatives are excluded, so $F$ is the whole chosen
subfield group. The identical argument applies to $E$.
\end{proof}

The proposition does not exclude two proper Suzuki subfield faces. In
particular neither equality of their fields nor descent of their joint
generated group has been proved. The chiral intersection property
$F\cap E=\gen b$ remains available, but is not a substitute for that
missing argument. Bounded trace probes recorded separately are optional
diagnostics and supply no uniform field-descent premise.
\end{document}
